\PassOptionsToPackage{table}{xcolor}
\usepackage{fontspec}
\setmainfont{Calibri}[Path=C:/Windows/Fonts/,UprightFont=calibri.ttf,BoldFont=calibrib.ttf,ItalicFont=calibrii.ttf,BoldItalicFont=calibriz.ttf]
\setmonofont{Consolas}[Path=C:/Windows/Fonts/,UprightFont=consola.ttf,BoldFont=consolab.ttf,ItalicFont=consolai.ttf,BoldItalicFont=consolaz.ttf,Scale=0.92]
\newfontfamily\symfont{Segoe UI Symbol}[Path=C:/Windows/Fonts/,UprightFont=seguisym.ttf]
\usepackage[most]{tcolorbox}
\usepackage{adjustbox,xcolor,titlesec,enumitem,needspace,fancyhdr,booktabs,array}
\renewcommand{\arraystretch}{1.3}
% Zebra-Streifen in Tabellen (sehr hell), Kopfzeile bleibt weiss
\rowcolors{2}{black!5}{white}
\setlength{\parindent}{0pt}
\setlength{\parskip}{2pt}
\setlist{nosep,leftmargin=1.2em}
\providecommand{\tightlist}{\setlength{\itemsep}{0pt}\setlength{\parskip}{0pt}}

% Farben
\definecolor{cUebung}{HTML}{0077A8}
\definecolor{cAufgabe}{HTML}{1F5FBF}
\definecolor{cLoesung}{HTML}{2E8B3D}
\definecolor{cHinweis}{HTML}{B8860B}
\definecolor{cWarn}{HTML}{C62828}
\definecolor{cErg}{HTML}{7B3FA0}
\definecolor{cBs}{HTML}{E07B00}

% Zeit-Badge (rechtsbündig, auch in Überschriften)
\newcommand{\zeitbadge}[1]{{\small\normalfont\color{black!65}{\symfont ⏱}\,#1}}
\newcommand{\zeit}[1]{\texorpdfstring{\hfill\zeitbadge{#1}}{}}

% Überschriften: ## = Phase (\section), ### = Schritt (\subsection)
\titleformat{\section}{\Large\bfseries\color{cUebung}}{}{0pt}{\needspace{10\baselineskip}}[{\vspace{-0.6ex}\color{cUebung}\titlerule[1.6pt]}]
\titlespacing*{\section}{0pt}{30pt}{8pt}
\titleformat{\subsection}{\large\bfseries\color{black!80}}{}{0pt}{\needspace{8\baselineskip}}
\titlespacing*{\subsection}{0pt}{20pt}{4pt}
\titleformat{\subsubsection}{\bfseries\color{black!70}}{}{0pt}{}

% Mermaid-Bilder
\newcommand{\mmdimg}[2][6.5cm]{\begin{center}\adjustbox{max width=1.0\linewidth,max height=#1}{\includegraphics[scale=0.4]{#2}}\end{center}}

% Boxen
\tcbset{boxstyle/.style={enhanced,breakable,left=5pt,right=5pt,top=3pt,bottom=3pt,boxrule=0.6pt,
  arc=2pt,fonttitle=\bfseries\small,before skip=5pt,after skip=5pt,toptitle=1pt,bottomtitle=1pt}}
\newtcolorbox{boxaufgabe}[1][Aufgabe]{boxstyle,colback=white,colframe=black!75,colbacktitle=cAufgabe!14,coltitle=cAufgabe,
  title={{\symfont ✎}\ #1}}
\newtcolorbox{boxloesung}[1][Musterlösung]{boxstyle,after skip=12pt,colback=white,colframe=black!75,colbacktitle=cLoesung!14,coltitle=cLoesung,
  title={{\symfont ✔}\ #1}}
\newtcolorbox{boxhinweis}[1][Hinweis]{boxstyle,colback=white,colframe=black!75,colbacktitle=cHinweis!14,coltitle=cHinweis,
  title={{\symfont ℹ}\ #1}}
\newtcolorbox{boxwarnung}[1][Nicht vergessen]{boxstyle,colback=white,colframe=black!75,colbacktitle=cWarn!14,coltitle=cWarn,boxrule=1.2pt,
  title={{\symfont ⚠}\ #1}}
\newtcolorbox{boxergaenzung}[1][Ergänzung (nicht auf der Webseite)]{boxstyle,colback=white,colframe=black!75,colbacktitle=cErg!14,coltitle=cErg,
  frame style={dashed},title={{\symfont ✚}\ #1}}
\newtcolorbox{boxbs}[1][Übergang ins betreute Selbststudium]{boxstyle,colback=white,colframe=black!75,colbacktitle=cBs!14,coltitle=cBs,
  title={{\symfont ➜}\ #1}}

% Kopf-/Fußzeile
\pagestyle{fancy}
\fancyhf{}
\renewcommand{\headrulewidth}{0pt}
\fancyhead[L]{\small\color{black!60}\leftmark}
\fancyfoot[R]{\small\color{black!60}Seite \thepage}
\fancyfoot[L]{\small\color{black!60}Dozentenskript -- nicht für Studierende}
\renewcommand{\sectionmark}[1]{}

% Code etwas kleiner
\AtBeginDocument{\ifcsname Highlighting\endcsname\RecustomVerbatimEnvironment{Highlighting}{Verbatim}{commandchars=\\\{\},fontsize=\small}\fi}
